\documentclass[12pt,a4paper]{report}
\usepackage[utf8]{vietnam}
\renewcommand{\contentsname}

\usepackage{amsmath}

\usepackage{amsfonts}
\usepackage{amssymb}
\usepackage{tikz}
\usetikzlibrary{arrows.meta, positioning}
\usepackage{listings}
\usepackage[dvipsnames]{xcolor}
%\usepackage[utf8]{inputenc}
%\usepackage[francais]{babel}
\usepackage{a4wide,amssymb,epsfig,latexsym,array,hhline,fancyhdr}
\usepackage[normalem]{ulem}
\usepackage{tocloft}
%\usepackage{soul}
\usepackage{caption}
\usepackage{float}
\renewcommand{\figurename}{}
\usepackage{wasysym}

\usepackage{enumitem}
\usepackage[makeroom]{cancel}
\usepackage{amsthm}
\usepackage{multicol,longtable,amscd}
\usepackage{diagbox}%Make diagonal lines in tables
\usepackage{booktabs}
\usepackage{alltt}
\usepackage[framemethod=tikz]{mdframed}
% For highlighting paragraph backgrounds
\usepackage{caption,subcaption}
\usepackage{listings}
\usepackage{color}
\usepackage{lastpage}
\usepackage[lined,boxed,commentsnumbered]{algorithm2e}
\usepackage{enumerate}
\usepackage{color}
\usepackage{graphicx}							
% Standard graphics package
\usepackage{array}
\usepackage{tabularx, caption}
\usepackage{multirow}
\usepackage{multicol}
\usepackage{rotating}
\usepackage{graphics}
\usepackage{geometry}
\usepackage{fontsize}
\usepackage{setspace}
\usepackage{epsfig}
\usepackage{indentfirst}
\usepackage{fontspec}
\usepackage{titlesec}
\titleformat{\chapter}[hang]
  {\normalfont\huge\bfseries}
  {}  % Để trống phần số
  {0pt}
  {}
\setmainfont{Times New Roman}
\usepackage{tikz}
\usetikzlibrary{arrows,snakes,backgrounds}
%\usepackage{pstcol} % PSTricks with the standard color package
\usepackage[normalem]{ulem}
\newtheorem{theorem}{{\bf Định lý}}
\newtheorem{property}{{\bf Tính chất}}
\newtheorem{proposition}{{\bf Mệnh đề}}
\newtheorem{corollary}[proposition]{{\bf Hệ quả}}
\newtheorem{lemma}[proposition]{{\bf Bổ đề}}
\theoremstyle{definition}
\newtheorem{exer}{Bài toán}
\def\thesislayout
{	% A4: 210 × 297
	\geometry
    {
		a4paper,
		total={160mm,240mm},  % fix over page
		left=30mm,
		top=30mm,
	}
}
\renewcommand{\cfttoctitlefont}{\hfill\Large\bfseries}
\renewcommand{\cftaftertoctitle}{\hfill}

% Thêm dấu chấm nối cho Section
\renewcommand{\cftsecleader}{\cftdotfill{\cftdotsep}}

% Đặt độ sâu mục lục
\setcounter{tocdepth}{3}
%Tao hf
\thesislayout
%\usepackage{fancyhdr}
\setlength{\headheight}{40pt}
\pagestyle{fancy}
\fancyhf{}%xoa thiet lap foot/head
%_________%
\fancyhead[L]{
 \begin{tabular}{rl}
    \begin{picture}(25,15)(0,0)
    \put(0,-8){\includegraphics[width=8mm, height=8mm]{hcmut.png}}
    %\put(0,-8){\epsfig{width=10mm,figure=hcmut.eps}}
   \end{picture}&
	%\includegraphics[width=8mm, height=8mm]{hcmut.png} & %
	\begin{tabular}{l}
		\textbf{\bf \ttfamily Trường Đại Học Bách Khoa - ĐHQG TP.HCM}\\
        \textbf{\bf \ttfamily Khoa Khoa học ứng dụng}
	\end{tabular} 	
 \end{tabular}
\fancyhead[R]{
	\begin{tabular}{l}
		\tiny \bf \\
		\tiny \bf 
	\end{tabular}  }
\fancyfoot[L]{\scriptsize \ttfamily Báo cáo bài tập lớn Vật Lí   - PH1003 \\ Niên khoá 2026 - 2027}
\fancyfoot[R]{\scriptsize \ttfamily Trang {\thepage}/\pageref{LastPage}}
  \fancyhead[R]{\thepage} % Hiển thị số trang góc phải (nếu muốn)
  \renewcommand{\headrulewidth}{0.4pt} % Đường kẻ ngang
}
\renewcommand{\headrulewidth}{0.3pt}
\renewcommand{\footrulewidth}{0.3pt}
% Cấu hình lại kiểu trang plain để nó hiển thị y hệt như fancy (hiện logo ở trang đầu chương)
\fancypagestyle{plain}  
{
    \fancyhf{} % Xóa sạch mặc định của plain
      % Gắn lại header logo Bách Khoa vào đây
      \fancyhead[L]{
        \begin{tabular}{rl}
          \begin{picture}(25,15)(0,0)
          \put(0,-8){\includegraphics[width=8mm, height=8mm]{hcmut.png}}
          \end{picture}&
          \begin{tabular}{l}
            \textbf{\bf \ttfamily Trường Đại Học Bách Khoa - ĐHQG TP.HCM}\\
            \textbf{\bf \ttfamily Khoa Khoa học ứng dụng}
          \end{tabular}     
        \end{tabular}
      }
      }
      \fancyhead[R]{\thepage} % Hiển thị số trang góc phải (nếu muốn)
      \renewcommand{\headrulewidth}{0.4pt} % Đường kẻ ngang
      \fancyfoot{} % clear all footer fields
        \fancyfoot[L]{\scriptsize \ttfamily GIẢI TÍCH 1 – MT1003 | LỚP L02 | NHÓM AC \\ Niên khoá 2026 - 2027}
      \fancyfoot[R]{\scriptsize \ttfamily Trang {\thepage}/\pageref{LastPage}}
      \renewcommand{\headrulewidth}{0.3pt}
      \renewcommand{\footrulewidth}{0.3pt}
      
}
\setcounter{secnumdepth}{4}
\setcounter{tocdepth}{3}
\makeatletter
\renewcommand{\thesection}{\arabic{section}}
\newcounter {subsubsubsection}[subsubsection]
\renewcommand\thesubsubsubsection{\thesubsubsection .\@alph\c@subsubsubsection}
\newcommand\subsubsubsection{\@startsection{subsubsubsection}{4}{\z@}%
                                     {-3.25ex\@plus -1ex \@minus -.2ex}%
                                     {1.5ex \@plus .2ex}%
                                     {\normalfont\normalsize\bfseries}}
\newcommand*\l@subsubsubsection{\@dottedtocline{3}{10.0em}{4.1em}}
\newcommand*{\subsubsubsectionmark}[1]{}
\makeatother
% --- ĐOẠN CODE LÀM ĐẬM SECTION TRONG MỤC LỤC ---
\makeatletter
\let\oldl@section\l@section
\renewcommand{\l@section}[2]{%
  \oldl@section{\bfseries#1}{\bfseries#2}%
}
\makeatother

\everymath{\color{black}}%make in-line maths symbols blue to read/check easily

\sloppy
\captionsetup[figure]{labelfont={small,bf},textfont={small,it},belowskip=-1pt,aboveskip=-9pt}
%space remove between caption, figure, and text
\captionsetup[table]{labelfont={small,bf},textfont={small,it},belowskip=-1pt,aboveskip=7pt}
%space remove between caption, table, and text

%\floatplacement{figure}{H}%forced here float placement automatically for figures
%\floatplacement{table}{H}%forced here float placement automatically for table
%the following settings (11 lines) are to remove white space before or after the figures and tables
%\setcounter{topnumber}{2}
%\setcounter{bottomnumber}{2}
%\setcounter{totalnumber}{4}
%\renewcommand{\topfraction}{0.85}
%\renewcommand{\bottomfraction}{0.85}
%\renewcommand{\textfraction}{0.15}
%\renewcommand{\floatpagefraction}{0.8}
%\renewcommand{\textfraction}{0.1}
\setlength{\floatsep}{5pt plus 2pt minus 2pt}
\setlength{\textfloatsep}{5pt plus 2pt minus 2pt}
\setlength{\intextsep}{10pt plus 2pt minus 2pt}

\thesislayout
\lstset{
language=C++,
    aboveskip=6mm,
    belowskip=-6mm,
    backgroundcolor=\color{white},   
    commentstyle=\color{codegreen},
    keywordstyle=\color{blue},
    numberstyle=\tiny\color{codegray},
    stringstyle=\color{codepurple},
    basicstyle=\ttfamily,
    breakatwhitespace=false,         
    breaklines=true,                 
    captionpos=b,                    
    keepspaces=true,                 
    numbers=left,                    
    numbersep=5pt,                  
    showspaces=false,                
    showstringspaces=false,
    showtabs=false,                  
    tabsize=2,
    frame=tb,
    framesep=8pt,
    rulecolor=\color{black},
    title=\lstname,
}

\usepackage{listings}

\usepackage{xcolor}

%New colors defined below
\definecolor{codegreen}{rgb}{0,0.6,0}
\definecolor{codegray}{rgb}{0.5,0.5,0.5}
\definecolor{codepurple}{rgb}{0.58,0,0.82}
\definecolor{backcolour}{rgb}{0.95,0.95,0.92}
\newfontfamily\listingunicodefont{DejaVu Sans Mono}

%Code listing style named "mystyle"
\lstdefinestyle{mystyle}{
  backgroundcolor=\color{backcolour}, commentstyle=\color{codegreen},
  keywordstyle=\color{magenta},
  numberstyle=\tiny\color{codegray},
  stringstyle=\color{codepurple},
  basicstyle=\ttfamily\footnotesize,
  breakatwhitespace=false,         
  breaklines=true,                 
  captionpos=b,                    
  keepspaces=true,                 
  numbers=left,                    
  numbersep=5pt,                  
  showspaces=false,                
  showstringspaces=false,
  showtabs=false,                  
  tabsize=2
}

\lstdefinestyle{pythonstyle}{
  language=Python,
  backgroundcolor=\color{backcolour},
  commentstyle=\color{codegreen},
  keywordstyle=\color{magenta}\bfseries,
  numberstyle=\tiny\color{codegray},
  stringstyle=\color{codepurple},
  basicstyle=\ttfamily\footnotesize,
  breaklines=true,
  captionpos=b,
  keepspaces=true,
  numbers=left,
  numbersep=6pt,
  showstringspaces=false,
  tabsize=4,
  frame=tb,
  framesep=6pt,
  rulecolor=\color{black},
  literate=
    {²}{{\listingunicodefont ²}}1 {³}{{\listingunicodefont ³}}1 {·}{{\listingunicodefont ·}}1 {À}{{\listingunicodefont À}}1
    {Á}{{\listingunicodefont Á}}1 {Â}{{\listingunicodefont Â}}1 {Ì}{{\listingunicodefont Ì}}1 {Í}{{\listingunicodefont Í}}1
    {Ô}{{\listingunicodefont Ô}}1 {Ù}{{\listingunicodefont Ù}}1 {à}{{\listingunicodefont à}}1 {á}{{\listingunicodefont á}}1
    {â}{{\listingunicodefont â}}1 {ã}{{\listingunicodefont ã}}1 {é}{{\listingunicodefont é}}1 {ê}{{\listingunicodefont ê}}1
    {ì}{{\listingunicodefont ì}}1 {í}{{\listingunicodefont í}}1 {ò}{{\listingunicodefont ò}}1 {ó}{{\listingunicodefont ó}}1
    {ô}{{\listingunicodefont ô}}1 {ú}{{\listingunicodefont ú}}1 {Đ}{{\listingunicodefont Đ}}1 {đ}{{\listingunicodefont đ}}1
    {Ĩ}{{\listingunicodefont Ĩ}}1 {ũ}{{\listingunicodefont ũ}}1 {Ơ}{{\listingunicodefont Ơ}}1 {ơ}{{\listingunicodefont ơ}}1
    {Ư}{{\listingunicodefont Ư}}1 {ư}{{\listingunicodefont ư}}1 {ạ}{{\listingunicodefont ạ}}1 {Ả}{{\listingunicodefont Ả}}1
    {ả}{{\listingunicodefont ả}}1 {ấ}{{\listingunicodefont ấ}}1 {ầ}{{\listingunicodefont ầ}}1 {ẫ}{{\listingunicodefont ẫ}}1
    {Ậ}{{\listingunicodefont Ậ}}1 {ậ}{{\listingunicodefont ậ}}1 {ằ}{{\listingunicodefont ằ}}1 {ặ}{{\listingunicodefont ặ}}1
    {Ẽ}{{\listingunicodefont Ẽ}}1 {ẽ}{{\listingunicodefont ẽ}}1 {Ế}{{\listingunicodefont Ế}}1 {ế}{{\listingunicodefont ế}}1
    {ể}{{\listingunicodefont ể}}1 {Ệ}{{\listingunicodefont Ệ}}1 {ệ}{{\listingunicodefont ệ}}1 {Ị}{{\listingunicodefont Ị}}1
    {ị}{{\listingunicodefont ị}}1 {ọ}{{\listingunicodefont ọ}}1 {ỏ}{{\listingunicodefont ỏ}}1 {ố}{{\listingunicodefont ố}}1
    {Ồ}{{\listingunicodefont Ồ}}1 {ồ}{{\listingunicodefont ồ}}1 {Ổ}{{\listingunicodefont Ổ}}1 {ổ}{{\listingunicodefont ổ}}1
    {ộ}{{\listingunicodefont ộ}}1 {ớ}{{\listingunicodefont ớ}}1 {Ờ}{{\listingunicodefont Ờ}}1 {ở}{{\listingunicodefont ở}}1
    {Ợ}{{\listingunicodefont Ợ}}1 {ợ}{{\listingunicodefont ợ}}1 {Ụ}{{\listingunicodefont Ụ}}1 {ụ}{{\listingunicodefont ụ}}1
    {Ứ}{{\listingunicodefont Ứ}}1 {ứ}{{\listingunicodefont ứ}}1 {Ừ}{{\listingunicodefont Ừ}}1 {ừ}{{\listingunicodefont ừ}}1
    {Ữ}{{\listingunicodefont Ữ}}1 {ữ}{{\listingunicodefont ữ}}1 {Ự}{{\listingunicodefont Ự}}1 {ự}{{\listingunicodefont ự}}1
    {₁}{{\listingunicodefont ₁}}1 {₂}{{\listingunicodefont ₂}}1 {→}{{\listingunicodefont →}}1 {√}{{\listingunicodefont √}}1
    {≠}{{\listingunicodefont ≠}}1
}

\begin{document}
\thispagestyle{empty}
\begin{center}
{\bfseries
    {\fontsize{16pt}{19pt}\selectfont
     ĐẠI HỌC QUỐC GIA TP. HỒ CHÍ MINH\\
     TRƯỜNG ĐẠI HỌC BÁCH KHOA\\
     KHOA KHOA HỌC ỨNG DỤNG\\
    }
     \vspace{1cm}
     \includegraphics[height=2cm]{hcmut.png}\\
     \vspace{2cm}
    {\fontsize{18pt}{22pt}\selectfont
    BÁO CÁO BÀI TẬP LỚN MÔN VẬT LÍ\\
    }
    \vspace{2cm}
    {\fontsize{18pt}{22pt}\selectfont
    BÀI 10: XÁC ĐỊNH LỰC TƯƠNG TÁC TĨNH ĐIỆN\\ GIỮA CÁC ĐIỆN TÍCH ĐIỂM TRONG CHÂN KHÔNG\par
    }
}
\vspace{2cm}
    {\fontsize{16pt}{19pt}\selectfont
\begin{flushleft}
    {\bfseries
    GVHD: Nguyễn Văn Khá\\
    \vspace{0.3cm}
    Tên Lớp/Tên Nhóm: Lớp L02/Nhóm 10\\
    \vspace{1cm}
    \centering
    Sinh viên thực hiện\\
    }
    \centering
\vspace{0.5cm}
\begin{tabular}{|l|l|}
\hline
\textbf{Họ tên} & \textbf{MSSV}\rule{0pt}{20pt}  \\ \hline
Nguyễn Thảo Vi& 2610353\rule{0pt}{20pt} \\ \hline
Nguyễn Hoàng Minh & 2613005\rule{0pt}{20pt} \\ \hline
Nguyễn Minh Nguyên  &2611431\rule{0pt}{20pt}\\ \hline

\end{tabular}
\end{flushleft} 
    }
\vfill
{\large TP. HCM, 2026\par}
\end{center}
\clearpage
\newpage
\tableofcontents
\section{MỞ ĐẦU}
\subsection{Cơ sở hình thành đề tài}
    Hiểu rõ các hiện tượng tự nhiên và kỹ thuật: Lực điện và điện trường đóng vai trò quan trọng trong nhiều hiện tượng tự nhiên (sét, hiện tượng nhiễm điện do cọ xát) cũng như các ứng dụng công nghệ hiện đại (máy lọc bụi tĩnh điện, công nghệ sơn tĩnh điện, máy in laser, các linh kiện vi điện tử).\\
    
    Giải quyết bài toán thực tế nhiều vật: Trong kỹ thuật và khoa học vật liệu, việc tính toán lực tương tác giữa nhiều hạt mang điện là bước cơ bản để mô phỏng cấu trúc nguyên tử, phân tử hoặc các hệ thống vi cơ điện tử (MEMS).\\
    
    Tuy nhiên, khi số lượng điện tích tăng lên (n ≥ 3), việc tính toán bằng tay trở nên phức tạp, tốn thời gian và dễ dẫn đến sai sót do các phép biến đổi đại số vectơ. Do đó, việc ứng dụng các công cụ tính toán ký hiệu (Symbolic Math) trên các ngôn ngữ lập trình như MATLAB hoặc Python giúp mô phỏng, tự động hóa và trực quan hóa các thành phần lực Fx, Fy một cách nhanh chóng và chính xác.\\
\subsection{Mục tiêu của đề tài}
    \begin{itemize}[leftmargin=1.2cm]
    \item Xây dựng thuật toán và chương trình (bằng MATLAB hoặc Python) ứng dụng phép toán hình thức (symbolic) để tính toán lực tương tác tĩnh điện.
    \end{itemize}
    \begin{itemize}
    \item Xử lý bài toán hệ n điện tích điểm trong mặt phẳng Oxy: Nhập tọa độ ($x_i, y_i$)
    và điện tích qiá trị, tự động xuất biểu thức và giá trị các thành phần lực $F_x, F_y$ tác dụng lên từng điện tích.
    \item Trực quan hóa hệ điện tích và vectơ lực tổng hợp tác dụng lên từng điện tích bằng đồ họa.
    \end{itemize}
\subsection{Ý nghĩa thực tiễn của đề tài}
    \begin{itemize}[leftmargin=1.2cm]
    \item Cung cấp công cụ học tập và kiểm tra kết quả nhanh chóng cho sinh viên khi làm bài tập Điện học.
    \end{itemize}
    \begin{itemize}[leftmargin=1.2cm]
    \item Rèn luyện kỹ năng lập trình mô phỏng vật lý, tư duy thuật toán và sử dụng thư viện đại số ký hiệu (SymPy / Symbolic Toolbox).
    \end{itemize}
\subsection{Phạm vi giới hạn của đề tài}
    \begin{itemize}[leftmargin=1.2cm]
    \item Hệ điện tích được xét là các điện tích điểm đặt cố định trong môi trường chân không (môi trường đẳng hướng, hằng số điện môi ε = 1).
    \end{itemize}
    \begin{itemize}[leftmargin=1.2cm]
    \item Không không gian khảo sát hạn chế trong mặt phẳng tọa độ hai chiều Oxy.
    \end{itemize}
\section{CƠ SỞ LÝ THUYẾT}

\subsection{Nêu những khái niệm, vấn đề lý thuyết có liên quan}

\subsubsection{Định luật Coulomb}
Lực tương tác tĩnh điện giữa hai điện tích điểm $q_1$ và $q_2$ đặt trong chân không cách nhau khoảng $r_{12}$ có độ lớn tỉ lệ thuận với tích độ lớn các điện tích và tỉ lệ nghịch với bình phương khoảng cách giữa chúng:
\[
\vec{F}_{12} = k \frac{|q_1 q_2|}{r_{12}^2} \vec{e}_r
\]
Trong đó $k = \frac{1}{4\pi\varepsilon_0} \approx 9 \times 10^9 \text{ N}\cdot\text{m}^2/\text{C}^2$.

\subsubsection{Nguyên lí chồng chất lực tĩnh điện}
Lực tổng hợp tĩnh điện do $n$ điện tích điểm tác dụng lên một điện tích $q_0$ bằng tổng các vectơ lực do từng điện tích $q_i$ tác dụng riêng rẽ lên $q_0$.

Biểu thức lực tổng hợp $\vec{F}_i$ tác dụng lên điện tích thứ $i$ và các thành phần theo hệ tọa độ $Oxy$. Trong đó khoảng cách giữa điện tích $q_i$ và $q_j$ được tính là:
\[
r_{ij} = \sqrt{(x_i - x_j)^2 + (y_i - y_j)^2}
\]

\subsection{Mô hình hóa hoặc hệ thống hóa vấn đề nghiên cứu}
Mô hình hóa hệ $n$ điện tích điểm trong mặt phẳng $Oxy$:

Xét điện tích $q_i$ tại vị trí $(x_i, y_i)$ và $q_j$ tại vị trí $(x_j, y_j)$. Vectơ khoảng cách $\vec{r}_{ij} = (x_j - x_i)\vec{i} + (y_j - y_i)\vec{j}$.

Khoảng cách:
\[
r_{ij} = \sqrt{(x_j - x_i)^2 + (y_j - y_i)^2}
\]

Vectơ lực do $q_i$ tác dụng lên $q_j$ chiếu lên trục $Ox$ và $Oy$ có dạng biểu thức đại số:
\begin{align*}
F_{ij,x} &= k \cdot \frac{q_i q_j (x_j - x_i)}{\left[(x_j - x_i)^2 + (y_j - y_i)^2\right]^{3/2}} \\
F_{ij,y} &= k \cdot \frac{q_i q_j (y_j - y_i)}{\left[(x_j - x_i)^2 + (y_j - y_i)^2\right]^{3/2}}
\end{align*}

\subsection{Xây dựng giả thiết}
\begin{itemize}[leftmargin=1.2cm]
    \item \textbf{Giả thiết 1:} Tương tác tĩnh điện tuân theo định luật III Newton, nghĩa là tổng lực tương tác tĩnh điện của toàn hệ kín phải bằng $0$: 
    \[
    \sum \vec{F}_j = \vec{0}
    \]
    \item \textbf{Giả thiết 2:} Phép toán hình thức (Symbolic) cho phép giữ nguyên dạng biểu thức chính xác tuyệt đối trước khi thế số, không gây ra lỗi tích lũy do làm tròn số lẻ.
\end{itemize}

\section{PHƯƠNG PHÁP NGHIÊN CỨU HOẶC THIẾT KẾ NGHIÊN CỨU}

\subsection{Chọn phương pháp}
Phương pháp thực hiện là mô phỏng số trên Python, kết hợp phép tính số thực (thư viện \texttt{math}), lập biểu thức ký hiệu (SymPy) và vẽ vectơ (Matplotlib). Đây là mô hình tính toán theo định luật Coulomb, không phải phép đo thực nghiệm trong phòng thí nghiệm.

\subsection{Nhu cầu thông tin và Nguồn thông tin}
\begin{itemize}[leftmargin=1.2cm]
    \item \textbf{Nhu cầu thông tin:} Công thức định luật Coulomb, thuật toán tính tổng đại số vectơ, tài liệu hướng dẫn thư viện SymPy/MATLAB.
    \item \textbf{Nguồn thông tin:} Giáo trình Vật lý đại cương A1 (Bộ môn Vật lý Ứng dụng -- ĐH Bách Khoa TP.HCM), tài liệu \textit{MATLAB Projects for Scientists and Engineers} (A. L. Garcia).
\end{itemize}

\subsection{Thiết kế mẫu (Thiết kế bài toán \& dữ liệu)}
Thiết kế mẫu dữ liệu kiểm thử gồm hệ 3 điện tích điểm đặt tại các đỉnh của một tam giác cân trong mặt phẳng $Oxy$:
\begin{itemize}[leftmargin=1.2cm]
    \item $q_1 = +2\ \mu\text{C}$ tại vị trí $(0;\, 0)\text{ m}$
    \item $q_2 = -4\ \mu\text{C}$ tại vị trí $(3;\, 0)\text{ m}$
    \item $q_3 = +3\ \mu\text{C}$ tại vị trí $(1{,}5;\, 2)\text{ m}$
\end{itemize}
Đây là cùng bộ dữ liệu mẫu được khai báo trong chương trình; các kiểm thử bổ sung bao gồm hai điện tích hút nhau, hai điện tích đẩy nhau, triệt tiêu đối xứng và điện tích trung hòa.

\subsection{Phương pháp và công cụ thu thập thông tin}
Mô hình giả sử các điện tích điểm đứng yên trong chân không, với $k=9\times10^9\ \text{N}\cdot\text{m}^2/\text{C}^2$; điện tích nhập bằng C và tọa độ bằng m. Các vị trí được giả định phân biệt nhau ($r_{ij}>0$), không xét ảnh hưởng của môi trường điện môi hay sự dịch chuyển điện tích. Vì vậy các kết quả là dự đoán số từ mô hình lý tưởng, không phải số liệu đo ngoài thực tế.

Chương trình nhận $n$, $q_i$ và $(x_i,y_i)$ qua bàn phím. Với từng điện tích đích $q_i$, vòng lặp xét mọi nguồn $q_j$ khác nó, tính $\Delta x=x_i-x_j$, $\Delta y=y_i-y_j$, $r^2=\Delta x^2+\Delta y^2$, rồi dùng
\[
\vec{F}_{j\to i}=k\,q_iq_j\frac{\Delta x\,\vec{i}+\Delta y\,\vec{j}}{(r^2)^{3/2}}.
\]
Chương trình cộng riêng các thành phần $F_x,F_y$, giữ các biểu thức theo hằng số ký hiệu $k$ bằng SymPy, sau đó tính độ lớn bằng \texttt{hypot} và hướng bằng \texttt{atan2}. Matplotlib vẽ các vị trí theo đúng tọa độ nhập vào; vectơ lực thành phần và hợp lực dùng cùng một hệ số thu phóng để hình dễ quan sát, nên độ dài mũi tên không biểu diễn trực tiếp số newton.

Các trích đoạn dưới đây được chia theo chức năng và lấy từ mã nguồn \texttt{../vatli1.py}.

\textbf{Phần 1 -- Khai báo hệ điện tích mẫu.} Mỗi phần tử lưu điện tích $q$ (C) và tọa độ $(x,y)$ (m); đây là bộ dữ liệu có thể chạy lại để đối chiếu kết quả trong báo cáo.
\lstinputlisting[style=pythonstyle,firstline=358,lastline=366,caption={Khai báo ba điện tích mẫu},label={lst:vatli1-data}]{../vatli1.py}

\textbf{Phần 2 -- Tính lực Coulomb do từng điện tích nguồn.} Với mỗi điện tích đang xét $q_i$, vòng lặp bỏ qua chính nó, xác định độ lệch tọa độ và khoảng cách, rồi tính các thành phần $F_x,F_y$ theo định luật Coulomb.
\lstinputlisting[style=pythonstyle,firstline=23,lastline=52,caption={Tính các thành phần lực theo định luật Coulomb},label={lst:vatli1-components}]{../vatli1.py}

\textbf{Phần 3 -- Chồng chất vectơ và lập kết quả.} Các thành phần từ mọi nguồn được cộng vào $F_{x,i}$ và $F_{y,i}$; chương trình dùng hai tổng này để tính độ lớn bằng \texttt{hypot}, đồng thời lưu biểu thức ký hiệu và lực thành phần để in hoặc vẽ.
\lstinputlisting[style=pythonstyle,firstline=53,lastline=92,caption={Cộng lực và lưu kết quả cho từng điện tích},label={lst:vatli1-result}]{../vatli1.py}

Để chạy kiểm tra từ thư mục \texttt{vatli1}, dùng lệnh \texttt{python tests/run\_cases.py}. Mỗi ca kiểm thử gồm cặp tệp cùng tên \texttt{.in}/\texttt{.out}: runner truyền nội dung \texttt{.in} vào stdin của bản sao chương trình CLI trong thư mục tạm, sau đó đối chiếu phần kết quả lực in ra với \texttt{.out}. Khi chạy chương trình bình thường, người dùng vẫn nhập dữ liệu tương tác qua bàn phím; bộ kiểm thử chỉ mô phỏng các lần nhập mẫu và không thay đổi luồng nhập liệu.

\subsection{Phương pháp phân tích dữ liệu hoặc kiểm định giả thiết}
\begin{itemize}[leftmargin=1.2cm]
    \item Đối chiếu thành phần lực và độ lớn với các giá trị chuẩn tính độc lập theo định luật Coulomb; kiểm thử gồm 5 tình huống trong \texttt{tests/cases}.
    \item Kiểm định định luật III Newton bằng cách tính tổng các thành phần $\sum_i F_{x,i}$ và $\sum_i F_{y,i}$ của toàn hệ.
\end{itemize}
\section{KẾT QUẢ}

\subsection{Mô tả dữ liệu}
Sau khi nhập mẫu dữ liệu 3 điện tích điểm vào chương trình, mã nguồn thực hiện xuất ra các biểu thức ký hiệu tổng quát và tính giá trị lực số dư tương ứng.\par

    \textbf{B}\textbf{ả}\textbf{ng 4.1. B}\textbf{ả}\textbf{ng d}\textbf{ữ}\textbf{ li}\textbf{ệ}\textbf{u đ}\textbf{ầ}\textbf{u vào và k}\textbf{ế}\textbf{t qu}\textbf{ả}\textbf{ thành ph}\textbf{ầ}\textbf{n l}\textbf{ự}\textbf{c tác d}\textbf{ụ}\textbf{ng lên t}\textbf{ừ}\textbf{ng đi}\textbf{ệ}\textbf{n tích:}\par
    \begin{center}\small
    \begin{longtable}{|p{3.50cm}|p{3.50cm}|p{3.50cm}|p{3.50cm}|}
    \hline
    \textbf{Điện tích} & \textbf{Thành phần $F_x$ (\text{N})} & \textbf{Thành phần $F_y$ (\text{N})} & \textbf{Độ lớn $F_{\text{tổng}}$ (\text{N})} \\ \hline
   $q_1$ ($+2\ \mu\text{C}$) & $+0{,}6000$ & $-0{,}4500$ & $0{,}7500$ \\ \hline
        $q_2$ ($-3\ \mu\text{C}$) & $-0{,}1680$ & $+0{,}3456$ & $0{,}3843$ \\ \hline
        $q_3$ ($+4\ \mu\text{C}$) & $-0{,}4320$ & $+0{,}1044$ & $0{,}4444$ \\ \hline
    \end{longtable}
    \end{center}

\subsection{Kiểm định giả thiết / Kết quả phân tích}
\begin{itemize}[leftmargin=1.2cm]
    \item Kiểm định tổng lực tác dụng lên toàn hệ:
\end{itemize}
\begin{align*}
\sum F_x &= (+0{,}6000) + (-0{,}1680) + (-0{,}4320) = 0{,}0000\text{ N} \\
\sum F_y &= (-0{,}4500) + (+0{,}3456) + (+0{,}1044) = 0{,}0000\text{ N}
\end{align*}

$\Rightarrow$ Giả thiết 1 hoàn toàn đúng, thỏa mãn định luật bảo toàn động lượng và định luật III Newton trong hệ kín.

\subsection{Bình luận về kết quả (so sánh, đánh giá)}
\begin{itemize}[leftmargin=1.2cm]
    \item \textbf{So sánh với phương pháp tính tay truyền thống:} Phương pháp Symbolic lập trình cho kết quả chính xác $100\%$, thời gian tính toán chỉ tốn vài millisecond so với 15--20 phút tính tay.
    \item \textbf{Đánh giá tính tổng quát:} Chương trình có thể mở rộng tính toán cho $n = 10, 50$ hay $100$ điện tích mà không cần thay đổi cấu trúc mã nguồn.
\end{itemize}
\section{KẾT LUẬN VÀ KIẾN NGHỊ}
\subsection{Khái quát hóa (generalization) về những điều tìm thấy}
    Đề tài đã hoàn thành xuất sắc các yêu cầu của Bài 10: Xây dựng thành công thuật toán tính toán lực tương tác tĩnh điện bằng phương pháp đại số ký hiệu (Symbolic Math) trên mặt phẳng Oxy.\par
\subsection{Ý nghĩa (implication)}
    Khẳng định hiệu quả vượt trội của việc kết hợp giữa lý thuyết Vật lý đại cương và các công cụ tin học hiện đại trong việc giải quyết bài toán lý thuyết phức tạp.\par
\subsection{Những hạn chế}
    \begin{itemize}[leftmargin=1.2cm]
    \item Bài toán mới xét trong không gian 2D ($Oxy$) và môi trường chân không liên tục.
    \end{itemize}
    \begin{itemize}[leftmargin=1.2cm]
    \item Chưa xét đến hiệu ứng phân cực trong các môi trường điện môi không đồng nhất.
    \end{itemize}
\subsection{Kiến nghị hướng nghiên cứu tiếp theo}
    \begin{itemize}[leftmargin=1.2cm]
    \item Phát triển chương trình lên không gian 3 chiều ($Oxyz$).
    \end{itemize}
    \begin{itemize}[leftmargin=1.2cm]
    \item Xây dựng giao diện người dùng GUI tương tác trực quan (Tkinter / PyQt).
    \end{itemize}
\section{TÀI LIỆU THAM KHẢO}
[1] Đại học Bách Khoa TPHCM, Sổ tay hướng dẫn thực hiện luận văn tốt nghiệp đại học, 2013.\par
[2] A. L. Garcia and C. Penland, MATLAB Projects for Scientists and Engineers, Prentice Hall, Upper Saddle River, NJ, 1996.\par
[3] Python Novice Course Materials \& Documentation: https://lucydot.github.io/python\_novice/\par
[4] Visualising Data with Python / MATLAB: https://lucydot.github.io/python\_novice/14-visualising-data/index.html\par
\end{document}